\paragraph{}
Our lives would not be possible without agriculture. We use goods like medicine, pesticides, and rubber everyday, all of which are made from plants. Plants are ubiquitous in our lives, and they fulfill some of our essential needs. Most importantly, plants feed people globally. But many of us still suffer from hunger. In 2024, it was estimated that between 638 and 720 million people did not have enough to eat \cite{StateFoodSecurity2025}. Plant dieases make this situation even worse by robbing us of food \cite{carvajal-yepesGlobalSurveillanceSystem2019, heHighThroughputRoboticPhenotyping2025, nazarovInfectiousPlantDiseases2020}. Therefore, they are a direct threat to our livlihoods.

\paragraph{}

Plant diseases make it difficult for crops to survive and stay healthy \cite{agriosINTRODUCTION2005}. They interfere with a plant's cells causing them to malfunction or die \cite{huberRelationshipNutritionPlant2012}. And each cell helps the plant carry out essential life-sustaining functions: conducting photosynthesis, and absorbing water and nutrients. So, when these cells are damaged, the plant's ability to obtain and use nutrients is limited. As a result, plant diseases stunt the growth, and decrease the yield and nutritional value of crops \cite{savaryCropLossesDue2012, vyskaTradeoffDiseaseResistance2016}. In some cases, the damage is so severe that farmers lose up to 98 of their crops \cite{nazarovInfectiousPlantDiseases2020}. However, the losses caused by plant diseases goes further than just the yield and quality of crops; they extend to rural communities, consumers, the environment, and governments \cite{savaryCropLossesDue2012}. Therefore, plant diseases threaten the health of crops, and this has disastrous far-reaching effects on our livlihoods. Clearly, we need solutions to deal with them.


%\paragraph{}
%In this section, the literature related to image-processing in tomato plant pathogen detection is reviewed. First a high-level overview of plant phenotyping is given, since plant phenotyping is foundational to image-processing methods in plant disease detection. Image-based phenotyping methods are then discussed, highlighting the advantages these methods have over traditional ones. Next, a classification of these solutions into three types is presented: traditional, machine learning, and deep learning solutions. And this classification is used to evaluate current image-processing solutions for plant pathogen detection across different crops, with emphasis on tomato plants. Finally, the main challenges and limitations of these methods are identified and possible solutions are presented.

%\paragraph{}
%This literature review provides a foundational understanding of image processing in plant phenotyping and directions for developing or investigating solutions for low-cost image-processing phenotyping tools in automated greenhouses. The focus is on the existing image processing solutions applied in low-cost automated greenhouse environments. We begin by examining the fundamental concepts of plant phenotyping and the various traits that are amenable to image processing for predicting plant health, growth, and stress. Specifically, we note that leaf phenotypic traits are a robust proxy indicator of physiological and metabolic factors affecting plant health \cite{zhangHighthroughputPhenotypingPlant2023}.

%\paragraph{}
%Thereafter, aspects of image processing in plant phenotyping are discussed. Notably, image processing and acquisition, as well as various image sensor technologies and segmentation methods, are reviewed in the context of plant phenotyping. Hyperspectral and RGB imaging are the most common acquisition methods, and segmentation techniques critically influence the overall accuracy, precision, and quality of phenotypical trait analysis \cite{choudhurySegmentationTechniquesChallenges2020, muller-linowPlantScreenMobile2019}.

%Lastly, the status, challenges, and potential innovations of low-cost automated image processing in greenhouse environments are discussed. In this vein, the literature indicates a lack of affordable and generalizable image processing tools that can analyse and process large amounts of image data on inexpensive, resource-constrained hardware and sensors \cite{gomesComputationalSustainabilityComputational2009, minerviniImageAnalysisNew2015}.